Speculative verification changes query batching and can alter the greedy output of a quantized deployment. We present a phase-scoped query-to-batch (QB) transformation that preserves query and mask indexing, shares K/V storage and makes the layout part of graph reuse while retaining batched projections and feed-forward computation. A six-session CUDA experiment spanning two dates reduces speculative request latency relative to explicit per-query attention (PQ) by 3.06\%, 6.17\% and 2.83\% for Qwen3, SmolLM2 and Llama, with complete reference-output agreement on the fixed panel. A separate 72-context replay finds bitwise PQ/QB correspondence at QK, softmax and AV. Vulkan state interventions expose component-substitution limits; a scoped backend port subsequently preserves the reference on real grouped-query verification requests. Controlled mapping experiments show that CUDA Graph replay and representation costs can remove apparent optimization opportunities, while a cheaper padded path can fail the specified reference. Two new four-input, fresh-context service panels qualify real two-query verification but establish no practical QB-over-PQ benefit: the point reductions are 1.07\% on RTX and 0.49\% on B570, with the RTX estimate unresolved under 5.21\% identical-control imbalance. Speculative service is 41.8\% lower than target-only service in the RTX panel and 2.62 times as costly in the B570 panel under their respective frozen policies. These results separate a reference-preserving representation, its implementation cost and the net value of drafting, and identify the state and backend conditions required for transfer.
